%=============================================================================
% GUIA DE MELHORES PRÁTICAS PARA COMENTÁRIOS EM LATEX
% Template DI PUC-Rio
%=============================================================================

% ESTRUTURA DE COMENTÁRIOS IMPLEMENTADA NO TEMPLATE:
%
% 1. CABEÇALHOS DE SEÇÕES PRINCIPAIS
%    - Linhas duplas com "=============" (77 caracteres)
%    - Título em MAIÚSCULAS centralizado
%    - Separação visual clara entre seções
%
% 2. SUBSEÇÕES E GRUPOS
%    - Linhas simples com "---" seguido de descrição
%    - Agrupamento lógico de pacotes relacionados
%    - Descrição funcional clara
%
% 3. COMENTÁRIOS INLINE
%    - Alinhamento à direita para melhor legibilidade
%    - Descrição concisa da função de cada linha
%    - Espaçamento consistente
%
% 4. EXEMPLOS E INSTRUÇÕES
%    - Blocos com "-----------------------------------------------------------------------------"
%    - Título descritivo (EXEMPLO: ...)
%    - Código indentado e bem formatado
%    - Instruções claras e detalhadas
%
% 5. ORGANIZAÇÃO HIERÁRQUICA
%    - Seções principais claramente delimitadas
%    - Subsecções logicamente agrupadas
%    - Ordem de importância e dependência respeitada

%=============================================================================
% CONVENÇÕES ADOTADAS
%=============================================================================

% LARGURA PADRÃO: 77 caracteres para quebras de linha
% INDENTAÇÃO: 2 ou 4 espaços para código comentado
% ALINHAMENTO: Comentários inline alinhados verticalmente quando possível
% IDIOMA: Português brasileiro para clareza e acessibilidade
% CAPITALIZAÇÃO: MAIÚSCULAS para títulos principais, Primeira Letra para subtítulos

%=============================================================================
% BENEFÍCIOS DESTA ESTRUTURA
%=============================================================================

% 1. LEGIBILIDADE APRIMORADA
%    - Navegação rápida entre seções
%    - Identificação imediata da função de cada bloco
%    - Estrutura visual clara e consistente
%
% 2. MANUTENIBILIDADE
%    - Modificações facilitadas pela organização lógica
%    - Adição de novos pacotes em locais apropriados
%    - Documentação autoexplicativa
%
% 3. COLABORAÇÃO
%    - Entendimento rápido por novos colaboradores
%    - Padrão consistente em todo o projeto
%    - Instruções claras para personalização
%
% 4. DEBUGGING
%    - Localização rápida de problemas
%    - Identificação de dependências entre pacotes
%    - Histórico de modificações mais claro

%=============================================================================
% COMO MANTER O PADRÃO
%=============================================================================

% 1. Sempre use comentários descritivos para novos pacotes
% 2. Mantenha o agrupamento lógico (matemática, figuras, formatação, etc.)
% 3. Atualize os exemplos quando necessário
% 4. Preserve o alinhamento e espaçamento
% 5. Use a estrutura hierárquica estabelecida
